\documentclass[aspectratio=169]{beamer}
\usepackage{tikz}
\usepackage{graphicx}
\usepackage{listings}
\usepackage{xcolor}
\usepackage{hyperref}
\usepackage[spanish,es-noquoting]{babel}
\usepackage[round]{natbib}

\lstset{
  basicstyle=\ttfamily\small,
  keywordstyle=\color{blue},
  stringstyle=\color{red},
  commentstyle=\color{gray},
  showstringspaces=false
}
% Minimal look
\usetheme{Warsaw}
%\setbeamertemplate{headline}
\setbeamertemplate{footline}{
  \leavevmode%
  \hbox{%
    % Left: Section name
    \begin{beamercolorbox}[wd=.33\paperwidth,ht=2.25ex,dp=1ex,left]{author in head/foot}%
      \hspace*{2ex}\insertsection
    \end{beamercolorbox}%
    % Center: Presentation title
    \begin{beamercolorbox}[wd=.34\paperwidth,ht=2.25ex,dp=1ex,center]{title in head/foot}%
      \insertshorttitle
    \end{beamercolorbox}%
    % Right: Slide numbers
    \begin{beamercolorbox}[wd=.33\paperwidth,ht=2.25ex,dp=1ex,right]{date in head/foot}%
      \insertframenumber{} / \inserttotalframenumber\hspace*{2ex}
    \end{beamercolorbox}%
  }%
  \vskip0pt%
}

\usecolortheme{default}
\usefonttheme{serif}
\graphicspath{{../}{./}}
\setbeamertemplate{caption}[numbered]

% Title info
\title[Module 6: Advanced Prompting]{Module 6: Advanced Prompting}
\subtitle{Few-Shot and Chain-of-Thought}
\author{Ricardo Reyes-Jimenez}
\institute{Capgemini}
\date{\today}



\begin{document}

% Title slide
\begin{frame}
  \titlepage
\end{frame}

% Reset background for normal slides
\usebackgroundtemplate{}



\begin{frame}[fragile]{Session Objective}


\textbf{By the end of this module, the Engagement Manager will be able to:}

1. Clone their personal or corporate writing style using \textbf{Few-Shot Prompting}.
2. Force the AI to reason sequentially using \textbf{Chain-of-Thought (CoT)}.
3. Apply these advanced techniques to generate complex SteerCo (Steering Committee) reports.

\end{frame}



\begin{frame}[fragile]{1. Cloning Your Style (Few-Shot Prompting)}


\textbf{The AI should write like you, not like a generic robot.}

\begin{itemize}
  \item \textbf{Zero-Shot:} No examples provided. Good for simple tasks (e.g., summarizing).
  \item \textbf{One-Shot:} Providing one example to define the format.
  \item \textbf{Few-Shot:} Providing 2 to 3 examples of input/output to teach the AI your tone, vocabulary, and preferred corporate phrasing.
\end{itemize}
> \textbf{Why it matters:} EMs spend hours editing AI drafts because they sound ``too enthusiastic'' or ``robotic''. Few-Shot eliminates this editing time \cite{fewshot2025}.

\end{frame}



\begin{frame}[fragile]{Exercise 1: Few-Shot Style Cloning}


\textbf{Context:} You need to announce a new Lead Architect joining the project. You want the email to sound exactly like your previous announcements.
\textbf{Task:} Use Copilot to write the email based on your past examples.

\begin{itemize}
  \item \textbf{Draft Prompt to Try:}
\end{itemize}
*"I need to announce that Sarah Connor is joining as our new Lead Architect on Monday. Write the email using the exact same tone and structure as these two previous examples I wrote:
[Example 1: 'Team, please welcome John...']
[Example 2: 'Hi everyone, starting tomorrow, David will be...']"*

\end{frame}



\begin{frame}[fragile]{2. Chain-of-Thought (CoT)}


\textbf{Do not ask for the conclusion in Step 1.}

\begin{itemize}
  \item \textbf{The Concept:} Forcing the AI to ``think step-by-step'' before giving the final answer.
  \item \textbf{The Benefit:} It drastically reduces hallucinations and logical errors in complex analytical tasks \cite{cot2026}.
  \item \textbf{The Method:} Divide the prompt into sequential instructions: \textit{Extract -> Analyze -> Format}.
\end{itemize}
\end{frame}



\begin{frame}[fragile]{Exercise 2: CoT for SteerCo Reports}


\textbf{Context:} You have 3 pages of messy, bulleted meeting notes from a chaotic weekly sync. You need a clean, 1-page executive summary for the SteerCo.
\textbf{Task:} Use CoT to force the AI to categorize the notes before summarizing them.

\begin{itemize}
  \item \textbf{Draft Prompt to Try:}
\end{itemize}
*"Process the following raw meeting notes [Paste Notes]. Do not give me the final summary immediately. Follow these steps sequentially:
\textbf{Step 1:} Extract all mentions of budget issues.
\textbf{Step 2:} Extract all mentions of timeline delays.
\textbf{Step 3:} Based ONLY on the findings in Step 1 and 2, generate a 3-bullet executive summary suitable for a C-level Steering Committee."*

\end{frame}



\begin{frame}[fragile]{Recap \& Next Session}


\textbf{Today we learned:}
\begin{itemize}
  \item How to use \textbf{Few-Shot} prompting to eliminate ``AI-sounding'' text and save editing time.
  \item How \textbf{Chain-of-Thought (CoT)} forces the AI to reason sequentially, yielding highly accurate executive summaries from messy data.
\end{itemize}
\textbf{Next Session:}
\textit{Module 7: MS Learn \& Information Retrieval.} We will explore how to use Copilot to query official Microsoft documentation and governance policies securely.
\end{frame}




\begin{frame}[allowframebreaks]{References}
  \bibliographystyle{apalike}
  \bibliography{references}
\end{frame}


\end{document}